\documentclass[11pt]{article}
\usepackage[margin=1in]{geometry}
\usepackage{amsmath,amssymb}
\usepackage{booktabs}
\usepackage{siunitx}
\usepackage{hyperref}
\usepackage{graphicx}
\usepackage{caption}
\usepackage{xcolor}

\sisetup{
  separate-uncertainty = true,
  per-mode = symbol,
}

\title{W4: COSY INFINITY Cross-Validation of FELsim Beamline Optimisation}
\author{Eremey Valetov}
\date{15 February 2026}

\begin{document}
\maketitle

\section{Objective}

Cross-validate FELsim's 11-stage sequential Nelder-Mead beamline optimisation
against COSY INFINITY's internal FIT command. Both codes solve the same problem:
matching Twiss parameters at the undulator entrance of the UH MkV FEL transport
line (118 elements, 23 quadrupoles) to the targets of
Weinberg, Fisher \& Li~\cite{weinberg2025}.

The goal is independent confirmation that the optimised quadrupole settings
achieve the required undulator matching --- not necessarily identical currents,
but equivalent final Twiss.


\section{Method}

\subsection{Translation to COSY FIT}

FELsim's 11-stage optimisation was translated into 11 sequential COSY FIT
blocks within a single FOX input file. Each FIT block:
\begin{itemize}
  \item Declares variables (quadrupole currents) and bounds,
  \item Re-executes the full lattice (via \texttt{LATTICE}),
  \item Evaluates a combined MSE objective from the Twiss/dispersion goals,
  \item Iterates until convergence ($\varepsilon = 10^{-8}$, $N_\mathrm{max} = 1000$).
\end{itemize}

Optimised variables persist between FIT blocks, so each subsequent stage
inherits the results of all prior stages.

\subsection{Key Implementation Details}

\paragraph{Combined MSE objective.}
COSY's FIT with individual objectives (\texttt{ENDFIT} with multiple
\texttt{OBJ}) resets to initial values when \emph{any} objective fails to
converge. A single combined MSE objective avoids this:
\begin{equation}
  \mathrm{OBJ}(1) = \frac{1}{N}\sum_{i=1}^{N} w_i\,(f_i - g_i)^2,
\end{equation}
where $f_i$ is the computed Twiss parameter, $g_i$ the target, and $w_i$ the
weight. The sum is evaluated without \texttt{CONS()}, preserving DA derivatives
for gradient-based optimisation (FIT algorithm~1).

\paragraph{Initial Twiss from beam distribution.}
FELsim computes Twiss from particle statistics of a Gaussian beam
($\sigma_x = \SI{0.8}{mm}$, $\varepsilon = \SI{0.101}{\pi.mm.mrad}$).
The initial beta function is $\beta_0 = \sigma_x^2/\varepsilon \approx
\SI{6.34}{m}$, not $\beta_0 = \SI{1}{m}$. Using the wrong initial Twiss
gives completely incorrect transport ($\alpha_x = -2.8$ instead of $\approx 0$
after Stage~1). This was the most critical implementation finding.

\paragraph{Dipole model.}
Hard-edge dipoles (\texttt{DIL} without \texttt{FC}) with edge angles from
the beamline specification. The \texttt{DIL} $\beta$ parameter is set to zero
(no fringe field integral), meaning edge kicks use $\tan(\epsilon)/\rho$
without the fringe field correction term. See Section~\ref{sec:physics} for
implications.


\section{Results}

\subsection{Beam and Target Parameters}

\begin{table}[h]
\centering
\caption{Beam parameters and undulator Twiss targets.}
\label{tab:params}
\begin{tabular}{@{}lS@{}}
\toprule
Parameter & {Value} \\
\midrule
Energy & \SI{40}{MeV} \\
Normalised emittance $\varepsilon_n$ & \SI{8}{\pi.mm.mrad} \\
Geometric emittance $\varepsilon$ & \SI{0.101}{\pi.mm.mrad} \\
Initial beam size $\sigma_{x,y}$ & \SI{0.8}{mm} \\
Initial $\beta_0$ & \SI{6.34}{m} \\
\midrule
Target $\beta_x$ & \SI{1.400}{m} \\
Target $\alpha_x$ & 0.470 \\
Target $\beta_y$ & \SI{0.242}{m} \\
Target $\alpha_y$ & 0.000 \\
\bottomrule
\end{tabular}
\end{table}

\noindent These parameters are identical for the S1 (2~ps) and S3 (0.5~ps)
configurations --- the transverse beam parameters and undulator Twiss targets
do not depend on bunch length, so the optimisation result is the same for both.


\subsection{COSY FIT Convergence}

The full 11-stage FIT converges in a single COSY execution.
Table~\ref{tab:twiss} compares the final Twiss at the undulator entrance.

\begin{table}[h]
\centering
\caption{Final Twiss at undulator entrance: COSY FR\,0, COSY FR\,1,
FELsim, and targets.}
\label{tab:twiss}
\begin{tabular}{@{}l SSSS @{}}
\toprule
{Parameter} & {Target} & {COSY FR\,0} & {COSY FR\,1} & {FELsim NM} \\
\midrule
$\beta_x$ (\si{m})  & 1.400  & 1.400  & 1.400  & 1.400  \\
$\alpha_x$          & 0.470  & 0.470  & 0.470  & 0.470  \\
$\beta_y$ (\si{m})  & 0.242  & 0.242  & 0.242  & 0.242  \\
$\alpha_y$          & 0.000  & 0.000  & 0.000  & 0.000  \\
\midrule
{MSE}               & {---}  & {4.50e-9} & {7.93e-8} & {2.89e-5} \\
\bottomrule
\end{tabular}
\end{table}

Both COSY runs match the targets to high precision. COSY's gradient-based
FIT (algorithm~1, using DA derivatives) converges $\sim 10^4\times$
more precisely than Nelder-Mead. FR\,0 and FR\,1 achieve comparable
MSE ($\sim 10^{-8}$--$10^{-9}$), confirming that fringe fields do not
degrade convergence when warm-starting is used.


\subsection{Quadrupole Current Comparison}

A detailed comparison of optimised quadrupole currents across FELsim,
COSY FR\,0, and COSY FR\,1 is given in Table~\ref{tab:currents_fr1}
(Section~\ref{sec:fr1}).

Stages~1--2 and 6--9 show excellent agreement across all three runs
($\Delta < \SI{0.1}{A}$). The principal discrepancy is Stage~5 (double
triplet, elements 33--43): COSY finds a local minimum with
\emph{negative} currents for elements 35 and 37, while FELsim
(constrained to $[0, 10]$~A) finds all-positive currents. A four-point
multistart over Stage~5 starting values confirmed that COSY's
gradient-based optimiser consistently converges to the negative-polarity
basin.

Stages~10 and 11 also differ from FELsim, partly because the different
Stage~5 outcome propagates through the downstream lattice, and partly
due to the dipole model difference discussed in Section~\ref{sec:physics}.


\subsection{Dispersion}

The COSY-optimised solution shows residual horizontal dispersion
$\eta_x \approx \SI{33}{mm}$ at the undulator entrance.
The Stage~11 objective includes dispersion zeroing at element~92
(5th chromaticity quad) with weight~0.5, but does not constrain
$\eta_x'$. Between element~92 and the undulator (element~117),
the non-zero $\eta_x'$ causes dispersion to accumulate through
the final triplet drifts.

FELsim shows comparable residual dispersion at element~92
(0.017 in optimizer units), confirming this is a lattice design
feature rather than a code discrepancy.


\subsection{Single-Dipole Transfer Matrix Comparison}
\label{sec:sandwich}

To quantify the dipole model difference, we compare the 6$\times$6
transfer matrix of individual DPW-DPH-DPW ``sandwiches'' between FELsim
and COSY under three fringe settings: \texttt{FR~0} (no fringe),
\texttt{FR~3} (third-order fringe), and \texttt{FR~3} with \texttt{FC}
(Enge coefficients, where available).

\begin{table}[h]
\centering
\caption{Vertical edge kick $M_{43}$ and Y-plane stability for representative
dipole triplets. FELsim uses a triangle-model fringe correction;
COSY uses the \texttt{DIL} element under different \texttt{FR} settings.}
\label{tab:sandwich}
\begin{tabular}{@{} l c SSSS @{}}
\toprule
{Dipole} & {$\theta$ ($^\circ$)} & {FELsim} & {COSY FR\,0} & {COSY FR\,3} & {COSY FR\,3+FC} \\
\midrule
\multicolumn{6}{@{}l}{\emph{$M_{43}$ (vertical kick, $\mathrm{rad/m}$)}} \\
Transport  & 1.5   & -7.42e-3 & -7.71e-3 & -6.78e-3 & {---} \\
Chicane    & 4.0   & -1.11e-1 & -1.21e-1 & -8.93e-2 & {---} \\
FC1        & 11.25 & -9.50e-1 & -1.04e+0 & -7.79e-1 & -2.27e-1 \\
\midrule
\multicolumn{6}{@{}l}{\emph{$|\mathrm{Tr}(M_y)/2|$ (stability, $<1$ required)}} \\
Transport  & 1.5   & 0.9997 & 0.9997 & 0.9997 & {---} \\
Chicane    & 4.0   & 0.9977 & 0.9975 & 0.9982 & {---} \\
FC1        & 11.25 & 0.9823 & 0.9805 & 0.9855 & 0.9957 \\
\bottomrule
\end{tabular}
\end{table}

Three findings emerge from Table~\ref{tab:sandwich}:
\begin{enumerate}
  \item FELsim's Y-kick lies between COSY FR\,0 and FR\,3 for every
    dipole type, but closer to FR\,0 (the hard-edge model).
  \item The discrepancy grows with bending angle: 3.9\% at 1.5$^\circ$,
    8.9\% at 4$^\circ$, and 10\% at 11.25$^\circ$ (FR\,0 vs FELsim).
  \item The \texttt{FC} Enge coefficients (available for the FC1/FC2
    chicane dipoles) produce much weaker Y-kicks (76\% less than FELsim),
    indicating that realistic fringe fields partially cancel the geometric
    edge focusing.
\end{enumerate}

The cumulative effect of all 14 dipole sandwiches in FELsim gives
$|\mathrm{Tr}(M_y)/2| = 1.203$ --- formally \emph{unstable} without
quadrupoles. The quads must therefore provide sufficient Y-plane focusing
to stabilise the beamline. When FELsim-optimised currents (calibrated to
FELsim's dipole model) are injected into COSY (with its different edge
kicks), the compensation breaks down, explaining the Y-plane instability
observed in cross-validation ($|\mathrm{Tr}(M_y)/2| = 29$).


\subsection{COSY FIT with First-Order Fringe Fields (FR\,1)}
\label{sec:fr1}

To assess the impact of dipole fringe fields on the optimisation, we
re-ran the full 11-stage COSY FIT with the \texttt{FR\,1} command enabled,
which includes first-order fringe field effects (Enge-model based
Runge-Kutta integration through the fringe region).

\paragraph{Cold-start failure.}
Running FR\,1 from the same default initial guesses as FR\,0 produces
poor convergence: MSE~$= 1.7 \times 10^{-1}$ at order~3 and
$5.7 \times 10^{-1}$ at order~1. The Y-plane becomes unstable
($|\mathrm{Tr}(M_y)/2| = 1.43$) because the changed dipole edge kicks
push the sequential FIT stages into incompatible local minima. The
gradient-based FIT (algorithm~1) cannot recover once early stages
converge to a wrong basin.

\paragraph{Warm-start from FR\,0.}
Using the converged FR\,0 quadrupole currents as initial guesses for
FR\,1 resolves the convergence problem completely. The optimiser needs
only small adjustments from the nearby FR\,0 solution:
\begin{center}
\begin{tabular}{@{}lSS@{}}
\toprule
{Run} & {MSE} & {$\Delta I_\mathrm{max}$ (A)} \\
\midrule
FR\,0 (baseline)         & 4.50e-9  & {---} \\
FR\,1 (cold start)       & 1.66e-1  & {---} \\
FR\,1 (warm from FR\,0)  & 7.93e-8  & 0.76 \\
\bottomrule
\end{tabular}
\end{center}
The warm-started FR\,1 run achieves MSE~$= 7.9 \times 10^{-8}$,
comparable to the FR\,0 baseline, with a maximum current change of
0.76~A. Stages~1--2 and 6--9 barely change; the main adjustments are
in the triplets (Stages~3, 5, 10--11).

Table~\ref{tab:currents_fr1} compares the quadrupole currents across
all three converged runs: FELsim, COSY FR\,0, and COSY FR\,1.

\begin{table}[h]
\centering
\caption{Optimised quadrupole currents: FELsim, COSY FR\,0, and
COSY FR\,1 (warm-started from FR\,0).}
\label{tab:currents_fr1}
\sisetup{table-format=2.4}
\begin{tabular}{@{} c l SSS c @{}}
\toprule
{Elem} & {Section} & {FELsim (A)} & {COSY FR\,0 (A)} & {COSY FR\,1 (A)} & {Stage} \\
\midrule
1   & First doublet     & 0.822  & 0.839  & 0.842  & 1 \\
3   & First doublet     & 1.043  & 1.052  & 1.057  & 1 \\
10  & Chromaticity 1    & 3.883  & 3.957  & 3.979  & 2 \\
16  & Triplet 1         & 2.240  & 1.993  & 1.994  & 3 \\
18  & Triplet 1         & 4.953  & 4.859  & 4.955  & 3 \\
20  & Triplet 1         & 3.426  & 3.499  & 3.599  & 3 \\
27  & Chromaticity 2    & 4.666  & 4.685  & 4.711  & 4 \\
33  & Double triplet    & 2.694  & 0.357  & 0.518  & 5 \\
35  & Double triplet    & 2.652  & $-$1.975 & $-$1.849 & 5 \\
37  & Double triplet    & 0.277  & $-$1.995 & $-$1.985 & 5 \\
50  & Chromaticity 3    & 4.674  & 4.685  & 4.696  & 6 \\
56  & IP doublet        & 3.122  & 3.168  & 3.218  & 7 \\
58  & IP doublet        & 3.313  & 3.321  & 3.361  & 7 \\
61  & Post-IP doublet   & 5.178  & 5.166  & 5.252  & 8 \\
63  & Post-IP doublet   & 4.043  & 4.011  & 4.080  & 8 \\
70  & Chromaticity 4    & 4.682  & 4.662  & 4.811  & 9 \\
76  & Triplet 2         & 3.934  & 3.869  & 3.953  & 10 \\
78  & Triplet 2         & 4.079  & 4.452  & 4.506  & 10 \\
80  & Triplet 2         & 0.014  & 0.608  & 0.592  & 10 \\
87  & Chromaticity 5    & 1.362  & 1.259  & 1.610  & 11 \\
93  & Final triplet     & 0.945  & 2.348  & 2.473  & 11 \\
95  & Final triplet     & 2.885  & 3.802  & 3.602  & 11 \\
97  & Final triplet     & 2.192  & 4.747  & 3.995  & 11 \\
\bottomrule
\end{tabular}
\end{table}

\paragraph{FR\,3 computational cost.}
Third-order fringe field computation (\texttt{FR\,3}) was attempted but
is prohibitively expensive: each lattice evaluation takes $>$15~minutes
due to Runge-Kutta integration through 14~dipole fringe regions, even at
first-order DA. A full 11-stage FIT (requiring hundreds of lattice
evaluations) would take days. The single-dipole comparison
(Table~\ref{tab:sandwich}) provides FR\,3 results for individual elements.


\section{Physics Model Differences}
\label{sec:physics}

\subsection{Dipole Edge Fringe Fields}

FELsim's \texttt{dipole\_wedge} class applies a simplified triangle-model
fringe field correction to the vertical edge kick:
\begin{equation}
  T_y = \tan\!\left(\epsilon - \frac{\ell_e}{6}\,\frac{1+\sin^2\!\epsilon}
        {\rho\,\cos\epsilon}\right),
\end{equation}
where $\ell_e$ is the physical DPW length, $\epsilon$ is the edge angle,
and $\rho$ is the bending radius. COSY's \texttt{DIL} with $\beta=0$
(our FR\,0 baseline) gives $T_y = \tan\epsilon$ without fringe correction,
while FR\,3 computes fringe fields to third order using the half-gap
parameter.

Each individual dipole shows a 4--20\% difference in the Y-kick
(Table~\ref{tab:sandwich}). Over 14 dipoles (28 edge kicks), this
accumulates to Y-plane instability when currents are exchanged between
codes.

COSY provides more physically accurate fringe field treatments:
the \texttt{FC} command with Enge function coefficients (measured values
are available for the FC1/FC2 dipoles), the \texttt{FR} command for
higher-order fringe fields, and the \texttt{MGE} command for 1D field map
integration. The \texttt{MGE} approach requires a correctly scaled
field map (currently unavailable for this beamline).

\subsection{Current Sign Handling}

FELsim's quadrupole model uses $k = |eGI/(m_e c \beta\gamma)|$,
making the transfer matrix independent of the sign of $I$.
COSY's \texttt{MQ} uses $B_\mathrm{pole} = \pm G I r$, where the sign
determines the focusing plane. Consequently, COSY's FIT can explore
negative-current solutions that reverse the focusing/defocusing planes,
while FELsim's Nelder-Mead (with bounds $[0, 10]$~A) is restricted to
the positive quadrant.

This explains why COSY consistently finds the negative-polarity Stage~5
solution: it is a genuine local minimum accessible to COSY but not to
FELsim's bounded optimiser. Both solutions satisfy the Stage~5 objectives.


\section{Summary}

\begin{enumerate}
  \item COSY INFINITY FIT successfully reproduces the 11-stage beamline
    optimisation in a single execution, achieving MSE~$\sim 10^{-8}$--$10^{-9}$
    for both FR\,0 and FR\,1 (warm-started).
    FELsim's Nelder-Mead achieves MSE~$= 2.9 \times 10^{-5}$.
  \item All three codes/settings match the undulator Twiss targets to
    high precision. The lower COSY MSE reflects the advantage of
    gradient-based (DA) optimisation over Nelder-Mead.
  \item Enabling first-order fringe fields (\texttt{FR\,1}) requires
    warm-starting from the FR\,0 solution. Cold-start with FR\,1 fails
    (MSE~$\sim 0.2$) because the changed dipole edge kicks push the
    sequential FIT stages into incompatible local minima.
  \item Third-order fringe fields (\texttt{FR\,3}) are prohibitively
    expensive: each lattice evaluation takes $>$15~minutes due to
    Runge-Kutta integration through 14 dipole fringe regions.
    Single-dipole comparisons (Table~\ref{tab:sandwich}) quantify
    the FR\,3 effect.
  \item The principal discrepancy is Stage~5, where COSY converges to
    negative-polarity currents while FELsim (constrained to $[0, 10]$~A)
    finds all-positive currents (Section~\ref{sec:physics}).
  \item Initial Twiss from the beam distribution ($\beta_0 = \sigma^2/\varepsilon
    \approx \SI{6.34}{m}$) is critical for correct transport computation.
  \item Direct current injection between codes fails due to the dipole edge
    fringe correction difference (5--8\% per edge, accumulating to Y-plane
    instability). The codes agree on final Twiss when each optimises
    within its own model.
  \item S1 and S3 produce identical transverse optimisation results.
\end{enumerate}

\paragraph{Next steps.}
\begin{itemize}
  \item Obtain the correct scaling coefficient for the chicane dipole field
    map to enable MGE-based simulation (the most physically accurate
    dipole model available in COSY).
  \item Add $\eta_x'$ or $\eta_x(117)$ objective to control undulator dispersion.
  \item Constrain Stage~5 variables to $I > 0$ (e.g., via $I^2$
    parameterisation) to find the positive-polarity COSY solution.
\end{itemize}


\begin{thebibliography}{9}
\bibitem{weinberg2025}
  M.~Weinberg, A.~Fisher, and R.~Li,
  ``Design of the UH Mark V Free Electron Laser,''
  arXiv:2510.14061v1, Table~I.
\end{thebibliography}

\end{document}
